% !TEX root = ../main.tex
%% Chapitre 10 — Polynômes et fractions rationnelles
%% Programme officiel MPSI/MP2I (BO spécial n°1 du 11/02/2021), premier semestre,
%% section « Polynômes et fractions rationnelles » : a) anneau K[X] ; b) division
%% euclidienne ; c) fonctions polynomiales et racines ; d) dérivation ;
%% e) arithmétique dans K[X] ; f) irréductibles de C[X] et R[X] ;
%% g) interpolation de Lagrange ; h) fractions rationnelles ; i) éléments simples.
%% Résultats vérifiés par verifications/chapitre_10.py.

\begin{ouverture}[polynomes]{10}{Polynômes et fractions rationnelles}{Partie 1 · Premier semestre}
  {Degré\enspace·\enspace Racines\enspace·\enspace Arithmétique de $\K[X]$\enspace·\enspace
   Factorisation\enspace·\enspace Lagrange\enspace·\enspace Éléments simples}
  \objectif{Calculer dans l'anneau $\K[X]$ : degrés, division euclidienne, algorithme de Horner.}
  \objectif{Relier racines, multiplicités, divisibilité et dérivées successives.}
  \objectif{Utiliser les relations de Viète et la formule de Taylor polynomiale.}
  \objectif{Transposer l'arithmétique de $\Z$ aux polynômes : PGCD, Bézout, Gauss.}
  \objectif{Factoriser dans $\Cx[X]$ et dans $\R[X]$, interpoler par Lagrange.}
  \objectif{Décomposer une fraction rationnelle en éléments simples et appliquer aux primitives.}
  \prerequis{Arithmétique de $\Z$ : division euclidienne, PGCD, Bézout, Gauss (chapitre 7).}
  \prerequis{Nombres complexes : racines de l'unité, second degré (chapitre 3).}
  \prerequis{Anneaux et corps (chapitre 8) ; dérivation (chapitre 4).}
\end{ouverture}

%% =====================================================================
%%  L'ESSENTIEL
%% =====================================================================
\essentiel

\essSection{1}{L'anneau $\K[X]$}

\begin{definition}[Polynômes]
Un \textbf{polynôme} à coefficients dans $\K$ ($=\R$ ou $\Cx$) s'écrit
$P=a_nX^{n}+\dots+a_1X+a_0$ ; si $a_n\neq0$, $n$ est le \textbf{degré}
$\deg P$, $a_n$ le \textbf{coefficient dominant} ($P$ est \textbf{unitaire}
si $a_n=1$). Par convention, $\deg0=-\infty$. L'ensemble $\K[X]$ est un
anneau commutatif \emph{intègre}, et $\K_n[X]$ désigne les polynômes de
degré au plus $n$. Pour tous $P,Q$ :
\[
  \deg(P+Q)\leq\max(\deg P,\deg Q),\qquad
  \deg(PQ)=\deg P+\deg Q,
\]
et pour la composée, $\deg(P\circ Q)=\deg P\times\deg Q$ (si $\deg Q\geq1$).
\end{definition}

\essSection{2}{Division euclidienne}

\begin{theoreme}
Pour tous $A,B\in\K[X]$ avec $B\neq0$, il existe un \emph{unique} couple
$(Q,R)$ tel que
\[
  A=BQ+R\qquad\text{et}\qquad\deg R<\deg B.
\]
On dit que $B$ \textbf{divise} $A$ lorsque $R=0$ ; deux polynômes se
divisant mutuellement sont \emph{associés} (proportionnels par un scalaire
non nul). La division se \emph{pose} comme pour les entiers, en éliminant à
chaque étape le terme de plus haut degré.
\end{theoreme}

\begin{savoirfaire}[L'algorithme de Horner]
Pour évaluer $P=a_nX^{n}+\dots+a_0$ en $c$, on calcule de proche en proche
$b_n=a_n$, puis $b_{k}=c\,b_{k+1}+a_{k}$ : le dernier terme est $P(c)$, et
les $b_k$ sont les coefficients du quotient de $P$ par $X-c$ — une
évaluation \emph{et} une division en $n$ multiplications.
\end{savoirfaire}

\essSection{3}{Racines, multiplicités, relations de Viète}

\begin{theoreme}[Racines et divisibilité]
$a$ est \textbf{racine} de $P$ si $P(a)=0$, ce qui équivaut à :
$(X-a)$ divise $P$. Un polynôme non nul de degré $n$ possède \emph{au plus
$n$ racines} (comptées avec multiplicité). Conséquences : un polynôme de
$\K_n[X]$ qui s'annule en $n+1$ points est nul, et deux polynômes de
$\K_n[X]$ qui coïncident en $n+1$ points sont égaux — un polynôme est
entièrement déterminé par sa fonction polynomiale.
\end{theoreme}

\begin{definition}[Multiplicité, polynôme scindé]
$a$ est racine de \textbf{multiplicité} $m$ de $P$ si $(X-a)^{m}$ divise
$P$ et $(X-a)^{m+1}$ ne le divise pas. $P$ est \textbf{scindé} sur $\K$
s'il s'écrit comme produit de facteurs de degré $1$ :
$P=a_n\prod_{i}(X-x_i)$.
\end{definition}

\begin{theoreme}[Relations de Viète]
Si $P=a_nX^{n}+\dots+a_0$ est scindé de racines $x_1,\dots,x_n$ :
\[
  x_1+x_2+\dots+x_n=-\frac{a_{n-1}}{a_n}
  \qquad\text{et}\qquad
  x_1x_2\cdots x_n=(-1)^{n}\,\frac{a_0}{a_n}.
\]
(Les relations intermédiaires se retrouvent en développant le produit.)
\end{theoreme}

\essSection{4}{Dérivation et formule de Taylor}

\begin{definition}
La \textbf{dérivée formelle} de $P=\sum a_kX^{k}$ est
$P'=\sum ka_kX^{k-1}$ ; elle est linéaire, vérifie $(PQ)'=P'Q+PQ'$ et la
formule de Leibniz. Pour $\K=\R$, elle coïncide avec la dérivée de la
fonction polynomiale.
\end{definition}

\begin{theoreme}[Taylor et multiplicités]
Pour tout $P\in\K_n[X]$ et tout $a\in\K$ :
\[
  P=\sum_{k=0}^{n}\frac{P^{(k)}(a)}{k!}\,(X-a)^{k}
  \qquad\text{(formule de Taylor polynomiale).}
\]
En conséquence, $a$ est racine de multiplicité $m$ de $P$ si et seulement
si
\[
  P(a)=P'(a)=\dots=P^{(m-1)}(a)=0
  \qquad\text{et}\qquad P^{(m)}(a)\neq0.
\]
En particulier : racine double $\iff$ racine commune de $P$ et $P'$.
\end{theoreme}

\essSection{5}{L'arithmétique de $\K[X]$}

Toute l'arithmétique du chapitre 7 se transpose, la division euclidienne
jouant le même rôle — seule la notion de « taille » change : le degré
remplace la valeur absolue.

\otIntertitre{Le dictionnaire $\Z\ \leftrightarrow\ \K[X]$}
{\renewcommand{\arraystretch}{1.45}\arrayrulecolor{otRule}%
\begin{tabularx}{\linewidth}{@{}>{\centering\arraybackslash}X>{\centering\arraybackslash}X@{}}
\toprule
{\otTete{Dans $\Z$}} & {\otTete{Dans $\K[X]$}} \\
\midrule
valeur absolue & degré \\
division euclidienne ($0\leq r<b$) & division euclidienne ($\deg R<\deg B$) \\
PGCD, algorithme d'Euclide & PGCD \emph{unitaire} $A\wedge B$, Euclide \\
Bézout, Gauss, PPCM & Bézout, Gauss, PPCM — à l'identique \\
nombres premiers & polynômes \emph{irréductibles} \\
décomposition en facteurs premiers & décomposition en facteurs irréductibles \\
\bottomrule
\end{tabularx}}

\vspace{2mm}
\begin{propriete}
Le PGCD de deux polynômes (non tous deux nuls) est l'unique diviseur commun
\emph{unitaire} de degré maximal, calculé par l'algorithme d'Euclide ; tout
diviseur commun le divise. Théorème de Bézout, lemme de Gauss, couples de
Bézout par l'algorithme étendu, PGCD d'une famille, polynômes premiers
entre eux (dans leur ensemble, deux à deux) : tous les énoncés du
chapitre 7 restent valables mot pour mot.
\end{propriete}

\essSection{6}{Factorisation dans $\Cx[X]$ et dans $\R[X]$}

\begin{theoreme}[d'Alembert-Gauss]
Tout polynôme non constant de $\Cx[X]$ possède au moins une racine
complexe (démonstration hors programme). Conséquences :
\begin{itemize}[label={\color{otNavy}\textbullet},leftmargin=1.4em,itemsep=0.6mm,topsep=0.8mm]
  \item les irréductibles de $\Cx[X]$ sont les polynômes de degré $1$, et
    tout polynôme de $\Cx[X]$ est \emph{scindé} ;
  \item dans $\Cx[X]$, la divisibilité se lit sur les racines et les
    multiplicités, et deux polynômes sont premiers entre eux si et
    seulement s'ils n'ont \emph{aucune racine commune} ;
  \item les irréductibles de $\R[X]$ sont les polynômes de degré $1$ et
    les polynômes de degré $2$ à discriminant strictement négatif ; deux
    racines complexes conjuguées d'un polynôme \emph{réel} ont la même
    multiplicité, et la décomposition réelle s'obtient en regroupant les
    racines conjuguées.
\end{itemize}
Exemple phare : $X^{n}-1=\displaystyle\prod_{\omega\in\mathbb{U}_n}(X-\omega)$
(chapitre 3).
\end{theoreme}

\essSection{7}{Interpolation de Lagrange}

\noindent
\begin{minipage}[c]{0.56\linewidth}
\begin{theoreme}
Soient $x_1,\dots,x_n\in\K$ \emph{distincts} et $y_1,\dots,y_n\in\K$. Il
existe un \emph{unique} $P\in\K_{n-1}[X]$ tel que $P(x_i)=y_i$ pour tout
$i$ :
\[
  P=\sum_{i=1}^{n}y_i\,L_i,\qquad
  L_i=\prod_{j\neq i}\frac{X-x_j}{x_i-x_j}.
\]
Les polynômes $Q$ (de degré quelconque) vérifiant $Q(x_i)=y_i$ pour tout
$i$ sont exactement les
\[
  Q=P+\Bigl(\prod_{i=1}^{n}(X-x_i)\Bigr)R,\qquad R\in\K[X].
\]
\end{theoreme}
\end{minipage}\hfill
\begin{minipage}[c]{0.4\linewidth}
\centering
\begin{tikzpicture}[scale=1.05,line cap=round]
  \draw[otMuted,line width=0.5pt,-{Stealth[length=1.6mm]}] (-0.4,0) -- (2.9,0);
  \draw[otMuted,line width=0.5pt,-{Stealth[length=1.6mm]}] (0,-0.3) -- (0,3.4);
  \draw[otNavy,line width=1.05pt]
    plot[domain=-0.18:2.52,samples=60] ({\x},{-1.5*\x*\x+3.5*\x+1});
  \foreach \px/\py in {0/1,1/3,2/2}{
    \fill[otGreenDark] (\px,\py) circle (1.6pt);
    \draw[otGreen!70!otMuted,densely dotted,line width=0.5pt] (\px,\py) -- (\px,0);}
  \node[below,font=\tiny,text=otMuted] at (0.06,-0.02) {$x_1$};
  \node[below,font=\tiny,text=otMuted] at (1,-0.02) {$x_2$};
  \node[below,font=\tiny,text=otMuted] at (2,-0.02) {$x_3$};
\end{tikzpicture}\\[0.5mm]
\legende{Trois points, une unique parabole : l'interpolation de Lagrange.}
\end{minipage}

\essSection{8}{Fractions rationnelles et éléments simples}

\begin{definition}
Le corps $\K(X)$ des \textbf{fractions rationnelles} est formé des
quotients $F=\dfrac{P}{Q}$ ($Q\neq0$), chacun admettant une \emph{forme
irréductible} ($P\wedge Q=1$, $Q$ unitaire). On pose
$\deg F=\deg P-\deg Q$ ; les racines de $P$ (resp. $Q$), dans la forme
irréductible, sont les \textbf{zéros} (resp. les \textbf{pôles}) de $F$,
avec leurs multiplicités. Si $\deg P\geq\deg Q$, la division euclidienne
fournit la \textbf{partie entière} : $F=E+\dfrac{R}{Q}$ avec $E\in\K[X]$
et $\deg R<\deg Q$.
\end{definition}

\begin{theoreme}[Décomposition en éléments simples]
Toute fraction rationnelle se décompose de manière unique en sa partie
entière plus une somme d'\emph{éléments simples} :
sur $\Cx$, des termes $\dfrac{a}{(X-\lambda)^{k}}$ ; sur $\R$, des termes
$\dfrac{a}{(X-\lambda)^{k}}$ et
$\dfrac{aX+b}{(X^{2}+pX+q)^{k}}$ ($\Delta<0$). En pratique (toute
technicité exclue) :
\begin{itemize}[label={\color{otNavy}\textbullet},leftmargin=1.4em,itemsep=0.6mm,topsep=0.8mm]
  \item \etape{Pôle simple} le coefficient de $\dfrac{1}{X-\lambda}$
    s'obtient en multipliant par $(X-\lambda)$ puis en évaluant en
    $\lambda$.
  \item \etape{Dérivée logarithmique} si
    $P=c\prod(X-\lambda_i)^{m_i}$, alors
    $\dfrac{P'}{P}=\displaystyle\sum_i\frac{m_i}{X-\lambda_i}$.
  \item \etape{Applications} calcul de primitives et de dérivées
    $k$-ièmes de fractions rationnelles.
\end{itemize}
\end{theoreme}

\begin{attention}[Pièges fréquents]
\begin{itemize}[label={\color{otOrange}\textbullet},leftmargin=1.4em,itemsep=0.6mm,topsep=0.5mm]
  \item $\deg(P+Q)$ peut chuter : les termes dominants peuvent se
    compenser.
  \item Le PGCD de polynômes est \emph{unitaire} par convention : ne pas
    oublier de normaliser.
  \item « Irréductible » dépend du corps : $X^{2}+1$ est irréductible sur
    $\R$, pas sur $\Cx$.
  \item Une racine de multiplicité $m$ de $P$ est racine de multiplicité
    $m-1$ de $P'$ — et non $m$.
  \item Avant de décomposer en éléments simples : forme irréductible
    \emph{puis} partie entière.
  \item Viète donne somme et produit \emph{pour un polynôme scindé} ; sur
    $\R$, vérifier d'abord que toutes les racines sont réelles — ou
    travailler dans $\Cx$.
\end{itemize}
\end{attention}

%% =====================================================================
%%  EXERCICES
%% =====================================================================
\exercices
\rubrique{Applications}

\begin{exercice}[1]{Degrés}
\begin{questions}
  \item Donner le degré et le coefficient dominant de
    $(2X^{3}-X+1)(X^{2}+3)$, de $(X^{2}+X)^{3}$ et de
    $(X^{3}+1)-(X^{3}-X)$.
  \item Donner un exemple où $\deg(P+Q)<\max(\deg P,\deg Q)$.
  \item Quel est le degré de $P\circ Q$ pour $\deg P=3$ et $\deg Q=2$ ?
\end{questions}
\end{exercice}

\begin{exercice}[1]{Division euclidienne}
\begin{questions}
  \item Poser la division euclidienne de $X^{3}+2X^{2}-X+1$ par
    $X^{2}+1$.
  \item Poser celle de $X^{4}-1$ par $X-2$.
  \item Quel est, sans calcul, le reste de la division d'un polynôme $P$
    par $X-a$ ? Vérifier sur la question 2.
\end{questions}
\end{exercice}

\begin{exercice}[1]{Racines évidentes et factorisation}
Soit $P=X^{3}-6X^{2}+11X-6$.
\begin{questions}
  \item Vérifier que $1$ est racine de $P$, et factoriser $P$ par $X-1$.
  \item Terminer la factorisation de $P$ en produit de facteurs de degré
    $1$.
  \item Contrôler le résultat à l'aide des relations de Viète (somme et
    produit des racines).
\end{questions}
\end{exercice}

\begin{exercice}[1]{Multiplicité d'une racine}
Soit $P=X^{4}-5X^{3}+9X^{2}-7X+2$.
\begin{questions}
  \item Calculer $P(1)$, $P'(1)$, $P''(1)$ et $P'''(1)$.
  \item En déduire la multiplicité de la racine $1$, puis factoriser $P$.
\end{questions}
\end{exercice}

\begin{exercice}[1]{Relations de Viète}
\begin{questions}
  \item Déterminer deux nombres dont la somme vaut $7$ et le produit
    $10$.
  \item Soient $x_1,x_2,x_3$ les racines (complexes) de
    $X^{3}-2X^{2}+3X-4$. Donner, sans les calculer, leur somme et leur
    produit.
  \item Calculer de même $x_1^{2}+x_2^{2}+x_3^{2}$.
    \emph{Développer $(x_1+x_2+x_3)^{2}$.}
\end{questions}
\end{exercice}

\begin{exercice}[1]{L'algorithme de Horner}
Soit $P=2X^{4}-3X^{3}+X-5$.
\begin{questions}
  \item Évaluer $P(2)$ par l'algorithme de Horner, en détaillant les
    étapes.
  \item En déduire, sans nouvelle division, le quotient et le reste de la
    division de $P$ par $X-2$.
  \item Combien de multiplications l'évaluation a-t-elle coûté ? Comparer
    avec l'évaluation naïve.
\end{questions}
\end{exercice}

\begin{exercice}[1]{Factoriser sur $\Cx$ puis sur $\R$}
Décomposer en facteurs irréductibles, dans $\Cx[X]$ puis dans $\R[X]$ :
\begin{questions}
  \item $X^{4}-1$ ;
  \item $X^{3}-1$ ;
  \item $X^{4}+X^{2}+1$.
    \emph{Écrire $X^{4}+X^{2}+1=(X^{2}+1)^{2}-X^{2}$.}
\end{questions}
\end{exercice}

\begin{exercice}[1]{PGCD de deux polynômes}
Soient $A=X^{3}-3X+2$ et $B=X^{2}-1$.
\begin{questions}
  \item Calculer $A\wedge B$ par l'algorithme d'Euclide.
  \item Retrouver le résultat en factorisant $A$ et $B$.
\end{questions}
\end{exercice}

\rubrique{Entraînement}

\begin{exercice}[2]{La formule de Taylor polynomiale}
\begin{questions}
  \item Écrire $X^{3}$ selon les puissances de $X-1$, à l'aide de la
    formule de Taylor en $1$.
  \item Vérifier le résultat en développant.
  \item En déduire sans calcul la valeur de
    $1+3h+3h^{2}+h^{3}$ pour tout $h$. Quel changement de variable
    venons-nous d'effectuer ?
\end{questions}
\end{exercice}

\begin{exercice}[2]{Racines doubles et paramètre}
Pour $\lambda\in\R$, on pose $P_\lambda=X^{3}-3X+\lambda$.
\begin{questions}
  \item Montrer que $P_\lambda$ admet une racine double si et seulement si
    cette racine est aussi racine de $P_\lambda'$.
  \item En déduire les valeurs de $\lambda$ pour lesquelles $P_\lambda$
    admet une racine double, et factoriser $P_\lambda$ dans chaque cas.
  \item Pour quelles valeurs de $\lambda$ le polynôme $P_\lambda$ a-t-il
    trois racines réelles distinctes ? \emph{Étudier la fonction
    $x\mapsto x^{3}-3x$.}
\end{questions}
\end{exercice}

\begin{exercice}[2]{La rigidité des polynômes}
\begin{questions}
  \item Montrer qu'un polynôme de $\K_n[X]$ qui s'annule en $n+1$ points
    distincts est le polynôme nul.
  \item Soit $P\in\R_n[X]$ tel que $P(k)=k$ pour $k=0,1,\dots,n$.
    Déterminer $P$.
  \item Soit $P\in\R[X]$ tel que $P(x+1)=P(x)$ pour tout réel $x$.
    Montrer que $P$ est constant.
\end{questions}
\end{exercice}

\begin{exercice}[2]{Polynômes réels et racines conjuguées}
\begin{questions}
  \item Montrer que si $P\in\R[X]$ et $z\in\Cx$ est racine de $P$, alors
    $\overline{z}$ est racine de $P$, avec la même multiplicité.
  \item On cherche un polynôme \emph{réel} unitaire de degré $4$ admettant
    $1+\mathrm{i}$ et $\sqrt2$ pour racines. Pourquoi $1-\mathrm{i}$ est-elle
    automatiquement racine, alors que $-\sqrt2$ ne l'est pas ? Déterminer un
    tel polynôme en choisissant $-\sqrt2$ comme quatrième racine.
  \item Pourquoi un polynôme réel de degré impair a-t-il toujours une
    racine réelle ? Donner deux arguments (chapitre 6 et chapitre 10).
\end{questions}
\end{exercice}

\begin{exercice}[2]{$X^{n}-1$ et les racines de l'unité}
\begin{questions}
  \item Rappeler la factorisation de $X^{n}-1$ dans $\Cx[X]$.
  \item En regroupant les racines conjuguées, factoriser $X^{5}-1$ dans
    $\R[X]$, en exprimant les coefficients à l'aide de
    $\cos\dfrac{2\pi}{5}$ et $\cos\dfrac{4\pi}{5}$.
  \item Donner les valeurs exactes de ces cosinus (chapitre 3), et écrire
    la factorisation entièrement explicite.
\end{questions}
\end{exercice}

\begin{exercice}[2]{Bézout chez les polynômes}
\begin{questions}
  \item Montrer que $A=X^{2}+1$ et $B=X-1$ sont premiers entre eux, de
    deux façons : par les racines (dans $\Cx$), puis par l'algorithme
    d'Euclide.
  \item Déterminer un couple $(U,V)$ tel que $UA+VB=1$.
  \item En déduire, pour tout $C\in\K[X]$, une solution de $UA+VB=C$.
\end{questions}
\end{exercice}

\begin{exercice}[2]{Interpolation de Lagrange}
\begin{questions}
  \item Déterminer l'unique polynôme $P\in\R_2[X]$ tel que $P(0)=1$,
    $P(1)=3$ et $P(2)=2$ : par les polynômes de base $L_i$, puis par
    identification.
  \item Décrire \emph{tous} les polynômes $Q\in\R[X]$ vérifiant ces trois
    conditions.
  \item Quel polynôme obtient-on en interpolant les points
    $(0,0)$, $(1,1)$, $(2,2)$, $(3,3)$ ? Commenter.
\end{questions}
\end{exercice}

\begin{exercice}[2]{Fractions : forme irréductible et partie entière}
Soit $F=\dfrac{X^{3}+1}{X^{2}-1}$.
\begin{questions}
  \item Mettre $F$ sous forme irréductible ; préciser zéros et pôles,
    avec leurs multiplicités.
  \item Déterminer la partie entière de $F$, et écrire $F$ sous la forme
    $E+\dfrac{R}{Q}$.
  \item En déduire que la courbe de la fonction associée admet une
    asymptote oblique en $\pm\infty$, que l'on précisera.
\end{questions}
\end{exercice}

\begin{exercice}[2]{Éléments simples : pôles simples}
Décomposer en éléments simples, sur $\R$ :
\begin{questions}
  \item $\dfrac{1}{X^{2}-1}$ ;
  \item $\dfrac{X+2}{X(X+1)}$ ;
  \item $\dfrac{1}{X(X^{2}+1)}$ ;
  \item En déduire une primitive de $x\mapsto\dfrac{1}{x^{2}-1}$ sur
    $\intOO{1}{+\infty}$.
\end{questions}
\end{exercice}

\begin{exercice}[2]{La dérivée logarithmique}
\begin{questions}
  \item Démontrer que si $P=c\,(X-\lambda_1)^{m_1}\cdots(X-\lambda_r)^{m_r}$,
    alors
    $\dfrac{P'}{P}=\displaystyle\sum_{i=1}^{r}\frac{m_i}{X-\lambda_i}$.
  \item Vérifier la formule pour $P=(X-1)^{2}(X+2)$.
  \item En déduire que les multiplicités se lisent sur $P'/P$ : comment
    reconnaître une racine simple ?
\end{questions}
\end{exercice}

\rubrique{Approfondissement}

\begin{exercice}[3]{Tchebychev : racines et factorisation}
On reprend les polynômes $T_n$ du chapitre 2
($T_0=1$, $T_1=X$, $T_{n+2}=2XT_{n+1}-T_n$), qui vérifient
$T_n(\cos t)=\cos(nt)$.
\begin{questions}
  \item Montrer que $\deg T_n=n$ et que le coefficient dominant de $T_n$
    vaut $2^{n-1}$ (pour $n\geq1$).
  \item Montrer que les $n$ réels
    $x_k=\cos\dfrac{(2k+1)\pi}{2n}$, $k=0,\dots,n-1$, sont des racines de
    $T_n$, deux à deux distinctes.
  \item En déduire que $T_n$ est scindé sur $\R$, à racines simples,
    toutes dans $\intOO{-1}{1}$, et écrire sa factorisation.
\end{questions}
\end{exercice}

\begin{exercice}[3]{Sommes de Newton}
Soient $x_1,x_2,x_3$ les racines complexes de $P=X^{3}-2X^{2}+3X-4$, et
$s_k=x_1^{k}+x_2^{k}+x_3^{k}$.
\begin{questions}
  \item Rappeler $s_1$ (exercice 5), puis $s_2$.
  \item En écrivant que chaque racine vérifie $x_i^{3}=2x_i^{2}-3x_i+4$,
    établir : $s_3=2s_2-3s_1+12$, et en déduire $s_3$.
  \item Généraliser : exprimer $s_{k+3}$ en fonction de $s_{k+2}$,
    $s_{k+1}$, $s_k$, et calculer $s_4$.
\end{questions}
\end{exercice}

\begin{exercice}[3]{Racines rationnelles}
\begin{questions}
  \item Soit $P=a_nX^{n}+\dots+a_0\in\Z[X]$ et $\dfrac{p}{q}$ une racine
    rationnelle de $P$, écrite sous forme irréductible. Montrer que
    $p\mid a_0$ et $q\mid a_n$.
    \emph{Multiplier $P\bigl(\frac pq\bigr)=0$ par $q^{n}$ et utiliser le
    lemme de Gauss.}
  \item Vérifier que $\sqrt2+\sqrt3$ est racine de $X^{4}-10X^{2}+1$.
  \item En déduire que $\sqrt2+\sqrt3$ est irrationnel — et comparer avec
    la démonstration du chapitre 5.
\end{questions}
\end{exercice}

\begin{exercice}[3]{L'erreur d'interpolation}
Soit $f$ de classe $\mathscr{C}^{3}$ sur $\intFF{a}{b}$,
$x_1<x_2<x_3$ dans $\intFF{a}{b}$, et $P\in\R_2[X]$ l'interpolant de $f$
en ces points. Soit $x\in\intFF{a}{b}$ distinct des $x_i$.
\begin{questions}
  \item On pose $K$ tel que
    $f(x)-P(x)=K\,(x-x_1)(x-x_2)(x-x_3)$, et
    $g(t)=f(t)-P(t)-K\,(t-x_1)(t-x_2)(t-x_3)$. Combien $g$ a-t-elle de
    zéros au moins ?
  \item En appliquant trois fois le théorème de Rolle (chapitre 6),
    montrer qu'il existe $\xi$ tel que $g'''(\xi)=0$, puis que
    $K=\dfrac{f'''(\xi)}{6}$.
  \item En déduire la majoration de l'erreur d'interpolation :
    \[
      |f(x)-P(x)|\leq\frac{M_3}{6}\,
      \bigl|(x-x_1)(x-x_2)(x-x_3)\bigr|,
      \qquad M_3=\max|f'''|.
    \]
\end{questions}
\end{exercice}

\begin{exercice}[3]{Une primitive par les éléments simples}
\begin{questions}
  \item Décomposer $\dfrac{1}{X^{3}+1}$ en éléments simples sur $\R$.
  \item En déduire, sur $\intOO{-1}{+\infty}$, une primitive de
    $x\mapsto\dfrac{1}{x^{3}+1}$ \emph{(il apparaîtra un logarithme et un
    arctangente)}.
  \item Contrôler le résultat en dérivant.
\end{questions}
\end{exercice}

\begin{exercice}[3]{Polynômes positifs et sommes de deux carrés}
On veut montrer que tout polynôme $P\in\R[X]$ vérifiant $P(x)\geq0$ pour
tout réel $x$ s'écrit $P=A^{2}+B^{2}$ avec $A,B\in\R[X]$.
\begin{questions}
  \item Vérifier l'identité, pour tous polynômes réels :
    \[
      (A^{2}+B^{2})(C^{2}+D^{2})=(AC-BD)^{2}+(AD+BC)^{2}.
    \]
  \item Montrer que, dans la factorisation réelle de $P$, les racines
    réelles ont une multiplicité \emph{paire}, et que le coefficient
    dominant est positif.
  \item En déduire que $P$ est un produit d'un carré $A_0^{2}$ et de
    facteurs $X^{2}+pX+q$ avec $\Delta<0$, chacun somme de deux carrés
    \emph{(forme canonique)}.
  \item Conclure à l'aide de la question 1, et écrire explicitement
    $X^{4}+3X^{2}+2$ comme somme de deux carrés de polynômes.
\end{questions}
\end{exercice}

% !TEX root = ../../main.tex
%% Défi du chapitre 10 — Le théorème de Gauss-Lucas
%% Résultats vérifiés par verifications/defis_06_10.py.

\rubrique{Défi}

\begin{defi}{Le théorème de Gauss-Lucas}{1 h 30}
Soit $P\in\Cx[X]$ de degré $n\geqslant2$, de coefficient dominant $\lambda$,
de racines distinctes $a_1,\dots,a_p$ de multiplicités $m_1,\dots,m_p$. Le
plan est identifié à $\Cx$. Le but est de montrer que les racines de $P'$
sont toujours situées « à l'intérieur » du nuage formé par les racines de
$P$ — plus précisément dans son \emph{enveloppe convexe}, l'ensemble des
barycentres à coefficients positifs des $a_k$.

\partieExo{A}{La dérivée logarithmique}
\begin{questions}
  \item Montrer l'égalité de fractions rationnelles
    $\dfrac{P'}{P}=\displaystyle\sum_{k=1}^{p}\frac{m_k}{X-a_k}$.
  \item Soit $z$ une racine de $P'$ qui n'est pas racine de $P$. Montrer
    que $\displaystyle\sum_{k=1}^{p}m_k\,\frac{z-a_k}{|z-a_k|^{2}}=0$.
    \emph{Conjuguer l'égalité obtenue en évaluant la question 1 en $z$.}
  \item En déduire le \emph{théorème de Gauss-Lucas} : toute racine de
    $P'$ s'écrit $\sum_k\theta_k a_k$ avec $\theta_k\geqslant0$ et
    $\sum_k\theta_k=1$.
\end{questions}

\medskip
\begin{center}
\begin{tikzpicture}[scale=1.25,line cap=round,line join=round]
  \draw[otMuted,line width=0.4pt,-{Stealth[length=1.5mm]}] (-2,0) -- (2.5,0);
  \draw[otMuted,line width=0.4pt,-{Stealth[length=1.5mm]}] (0,-1.7) -- (0,1.9);
  \fill[otNavy!10] (2,0) -- (1,1.5) -- (-1.5,1) -- (-1,-1.2) -- (0.8,-1.4) -- cycle;
  \draw[otNavy,line width=0.7pt,dash pattern=on 2.5pt off 1.5pt] (2,0) -- (1,1.5) -- (-1.5,1) -- (-1,-1.2) -- (0.8,-1.4) -- cycle;
  \foreach \p in {(2,0),(1,1.5),(-1.5,1),(-1,-1.2),(0.8,-1.4)}{\fill[otNavy] \p circle (2.2pt);}
  \foreach \p in {(-0.832,0.505),(-0.216,-0.876),(1.011,0.617),(1.077,-0.325)}{%
    \filldraw[fill=otGreen,draw=otGreenDark,line width=0.5pt] \p circle (2pt);}
  \draw[otInk,line width=0.6pt] (0.26,-0.02) ++(-0.07,-0.07) -- ++(0.14,0.14) ++(-0.14,0) -- ++(0.14,-0.14);
  \node[font=\scriptsize,text=otNavy,anchor=south west] at (1,1.5) {racines de $P$};
  \draw[otGreenDark,line width=0.4pt] (-0.216,-0.95) -- (-0.45,-1.62);
  \node[font=\scriptsize,text=otGreenDark,anchor=north] at (-0.45,-1.62) {racines de $P'$};
  \node[font=\scriptsize,text=otInk,anchor=west] at (0.33,0.08) {$g$};
\end{tikzpicture}\\[0.5mm]
\legende{Un polynôme de degré $5$ : ses racines (en bleu) délimitent un
pentagone convexe, à l'intérieur duquel se trouvent les quatre racines de
$P'$ (en vert) ; $g$ est la moyenne commune des deux familles (question 7).}
\end{center}

\partieExo{B}{Exemples}
\begin{questions}[resume]
  \item Déterminer les racines de $P'$ lorsque $P=X^{n}-1$, et
    interpréter.
  \item Soit $P=(X-1)(X+1)(X-2\mathrm{i})$. Calculer $P'$ et ses racines,
    puis écrire chacune d'elles comme barycentre à coefficients positifs de
    $1$, $-1$ et $2\mathrm{i}$. Faire une figure.
  \item On suppose $P\in\R[X]$ scindé sur $\R$, de racines
    $a_1<\dots<a_p$. Retrouver, à l'aide du théorème de Rolle (chapitre 6)
    et des multiplicités, que $P'$ est scindé sur $\R$ avec toutes ses
    racines dans $\intFF{a_1}{a_p}$.
\end{questions}

\partieExo{C}{Conséquences}
\begin{questions}[resume]
  \item À l'aide des relations de Viète, montrer que la moyenne des
    racines de $P'$ (comptées avec multiplicité) est égale à celle des
    racines de $P$.
  \item Montrer que si toutes les racines de $P$ sont dans un disque
    fermé $D$, il en est de même des racines de $P',P'',\dots,P^{(n-1)}$.
    Que vaut l'unique racine de $P^{(n-1)}$ ?
\end{questions}
\end{defi}

